\begin{textsample}{POS Dim 3 – 2009\_c6 – Score -3.00 – 2009\_c6/enade\_2009\_proc\_ger\_1f5954ceea.txt}  \label{ex:f3_pos_242}
Capítulo 3 Análise Técnica da Prova Este capítulo tem por objetivo apresentar o desempenho dos estudantes de TECNOLOGIA EM PROCESSOS GERENCIAIS no ENADE/2009.
Para isso, foram calculadas as estatísticas básicas da prova como um todo, bem como as estatísticas das questões relacionadas à Formação Geral e ao Componente Específico.
Nas Tabelas, são apresentadas as seguintes estatísticas : total da população e de presentes, média, erro-padrão da média, desvio-padrão, mediana, nota máxima, nota mínima e coeficiente de assimetria ( ver definição em 1.4.3 ).
Nas estatísticas das questões discursivas, foram detalhadas também as notas nulas : em branco, por protesto ou por insuficiência.
As estatísticas apresentadas neste capítulo contemplam, separadamente, os ingressantes, os concluintes e o total de estudantes, e foram calculadas tendo -se em vista as seguintes agregações : ( a ) Região Geográfica e País e ( b ) a Categoria Administrativa e a Organização Acadêmica.
As notas brutas das questões discursiva em cada uma das partes, Formação Geral/Componente Específico, foram calculadas a partir da média simples das notas de cada uma das questões correspondentes.
A primeira questão discursiva da Formação Geral foi desconsiderada e, portanto, a nota bruta referente às questões discursivas de Formação Geral foi igual à nota da questão mantida, a segunda.
Na componente específica, a nota bruta correspondeu à média das notas nas três questões.
As notas brutas das questões objetivas de cada uma das partes, Formação Geral/Componente Específico, foram calculadas a partir dos acertos dos estudantes considerando -se uma regra de três para a conversão para a escala de 0 a 100 descontando -se aquelas questões que foram desconsideradas pela comissão ou que não passaram pelo teste do índice de discriminação ponto-bisserial ( ver seção 1.4.6 para a descrição do índice ).
Questões com índices de discriminação inferior a 0,2 foram excluídas do cômputo da nota do estudante.
As notas brutas de Formação Geral e no Componente Específico para cada um dos estudantes foram calculadas como médias ponderadas das suas notas nas questões objetivas e discursivas, ponderação diferente para cada parte : C\_FG = C\_FG^OBJ × 0,6 + C\_FG^DIS × 0,4 C\_CE = C\_CE^OBJ × 0,85 + C\_CE^DIS × 0,15 onde C\_FG é a nota bruta da Formação Geral, C\_FG\_OBJ é a nota bruta nas questões objetivas da Formação Geral, C\_FG\_DIS é a nota bruta nas questões discursivas da Formação Geral, C\_CE é a nota bruta do Componente Específico, C\_CE\_OBJ é a nota bruta nas questões objetivas do Componente Específico, e C\_CE\_DIS é a nota bruta nas questões discursivas do Componente Específico.
A nota bruta da prova, C\_GER, por sua vez, é também uma média ponderada da nota de cada uma das partes, Formação Geral/Componente Específico : C\_GER = 0,25 × C\_FG + 0,75 × C\_CE. 3.1 Estatísticas Básicas da Prova 3.1.1 Estatísticas Básicas Gerais A Tabela 3.1 apresenta as estatísticas básicas da prova por grupo de estudantes.
O total de estudantes inscritos para a prova foi de 29.633.
Destes, 64,0 \% compareceram, sendo que a presença foi menor entre concluintes ( 62,6 \% ) do que entre ingressantes ( 65,0 \% ).
A média geral da prova foi de 38,24, sendo que os ingressantes obtiveram média mais baixa, de 37,43, que os concluintes, cuja média foi 39,48.
O desvio-padrão geral foi de 14,1, sendo o do grupo de ingressantes menor ( 13,8 ), do que o do grupo dos concluintes ( 14,4 ), aqueles possuíram uma distribuição mais homogênea nas notas.
A nota máxima foi 91,3, obtida por pelo menos um concluinte ao passo que a maior nota obtida por um ingressante foi 87,7.
Os coeficientes de assimetria para os dois grupos considerados, o de concluintes e o de ingressantes, bem como para a população como um todo foram todos positivos, indicando uma maior concentração à esquerda da distribuição e um maior espalhamento à direita.
Uma comparação entre o desempenho de ingressantes e concluintes pode ser realizada a partir dos dados apresentados no Gráfico 3.1.
Cumpre destacar que as notas dos concluintes, de um modo geral, foram mais elevadas do que as dos ingressantes, cuja distribuição de notas está à esquerda dos primeiros.
A impressão de melhor desempenho é reforçada pelas notas medianas dos dois grupos : 38,5 entre os concluintes e 36,6 entre os ingressantes.
Nas faixas de notas até 40 pontos, o percentual de ingressantes é maior que o dos concluintes, situação esta que se inverte para as faixas de notas superiores a 40 pontos.
Constatou -se que aproximadamente 18,1 \% dos ingressantes e 22,5 \% dos concluintes situam -se nas faixas de notas superiores a 50,0 pontos.
O Gráfico 3.2 apresenta informações referentes ao desempenho geral de ingressantes e concluintes, comparando os resultados considerando a Região, Categoria Administrativa e Organização Acadêmica.
Levando -se em conta as notas médias dos estudantes em cada Região, observou -se nos dois grupos, uma variação expressiva entre as Regiões.
As notas médias mais elevadas entre os concluintes ocorreram nas Regiões Sudeste ( 41,51 ) e Centro-Oeste ( 39,71 ), ambas acima da média nacional ( 39,48 ).
Em relação aos ingressantes, as médias mais altas foram obtidas na Região Sudeste ( 38,38 ) nas demais Regiões a média foi inferior à média nacional ( 37,42 ).
Na Região Norte foi obtido à média de notas mais baixa entre os concluintes ( 37,41 ) e na Região Centro-Oeste entre os ingressantes ( 34,73 ).
Cursos nesta área foram oferecidos em instituições federais, estaduais, municipais e particulares.
Levando -se em conta os agrupamentos dos estudantes em categorias administrativas, observa -se que a pontuação média mais elevada entre os concluintes e ingressantes foi encontrada em instituições de origem estadual ( respectivamente 47,56 e 45,32 ) – ambas acima da média nacional.
Tendo como foco as notas médias dos estudantes agrupados por Organização Acadêmica, notamos que Institutos Federais de Educação, Ciência e Tecnologia apresentam as maiores médias entre os concluintes ( 45,96 ) e Centro Universitário apresentam as maiores médias entre os ingressantes ( 38,87 ), ao passo que as piores notas médias foram observadas nas Universidades para concluintes ( 38,48 ) e Institutos Federais de Educação, Ciência e Tecnologia para os ingressantes ( 35,15 ). 3.1.2 Estatísticas Básicas de Formação Geral A Tabela 3.2 apresenta as estatísticas básicas de Formação Geral por grupo de estudantes.
O total de estudantes inscritos para a prova foi de 29.633.
Destes, 64,0 \% compareceram, sendo que a presença foi menor entre concluintes ( 62,6 \% ) do que entre ingressantes ( 65,0 \% ).
A média geral da parte da prova de Formação Geral foi de 41,02, sendo que os ingressantes obtiveram média mais baixa, 40,35, que os concluintes, cuja média foi 42,05.
O desvio-padrão geral foi de 20,4, sendo o do grupo de ingressantes ( 20,4 ) igual ao grupo dos concluintes, indicando igualdade no grau de homogeneidade das notas.
A nota máxima foi 100,0, obtida por estudantes de ambos os grupos.
Para este componente, assim como para a nota geral, os coeficientes de assimetria são positivos para todos os grupos.
O Gráfico 3.3 permite a comparação entre concluintes e ingressantes relativamente ao desempenho na parte da prova que aborda a Formação Geral.
Nesta parte, se excluídas as provas com nota zero, tal como na prova como um todo, também os concluintes apresentaram desempenho superior, embora seja perceptível um maior equilíbrio.
A mediana dos concluintes foi superior à dos ingressantes, o mesmo acontecendo com a média.
Dos ingressantes, aproximadamente 31,7 \% alcançaram nota superior a 50 ; dentre os concluintes, o percentual dos estudantes nesta faixa de notas foi muito semelhante e igual a 34,7 \%.
No Gráfico 3.4, são apresentadas as informações referentes ao desempenho dos ingressantes e dos concluintes na parte de Formação Geral, em diferentes agregações como Região, Categoria Administrativa e Organização Acadêmica.
Considerando -se as notas médias dos estudantes segundo cada Região, observou -se que, tanto no grupo dos ingressantes quanto no grupo dos concluintes, houve diferenças marcantes.
As médias variam, entre os concluintes, de 33,58 na Região Norte a 44,08 na Região Sudeste.
Entre os ingressantes, os valores foram de 32,67 na Região Centro-Oeste a 41,87 na Região Sudeste.
A maior dispersão de notas dos estudantes de uma dada Região foi encontrada no Norte tanto para os concluintes quanto para os ingressantes.
Nas Regiões Nordeste, Sudeste, Sul e Centro-Oeste os concluintes apresentaram médias maiores do que os ingressantes, a situação inversa aconteceu no Norte.
Todas as Regiões apresentaram concluintes.
Tendo como foco as categorias administrativas, destacou -se que todas as categorias estavam representadas.
Sendo que as notas médias mais elevadas tanto para ingressantes ( 53,05 ) quanto para concluintes ( 53,40 ) foram nas instituições estaduais.
As notas médias mais baixas para ingressantes ( 36,52 ) ocorrem nas instituições federais e para concluintes ( 40,02 ) nas instituições municipais.
Considerando -se o tipo de Organização Acadêmica, os Institutos Federais de Educação, Ciência e Tecnologia apresentaram os maiores valores para concluintes ( 46,43 ).
E os Centros Universitários apresentaram os maiores valores para ingressantes ( 42,75 ). 3.1.3 Estatísticas Básicas do Componente Específico A Tabela 3.3 apresenta as estatísticas básicas referentes ao Componente Específico da prova, tendo sua média geral, de 37,36, inferior ao desempenho da prova de TECNOLOGIA EM PROCESSOS GERENCIAIS como um todo cuja média foi 38,24, como apresentado na Tabela 3.1, que por sua vez é inferior ao desempenho na Formação Geral, que teve média igual a 41,02, conforme mostrado na Tabela 3.2.
Os concluintes apresentaram um desempenho médio ( 38,68 ) superior ao dos ingressantes ( 36,50 ), havendo ainda indícios de maior homogeneidade entre as notas médias dos ingressantes do que entre as dos concluintes, como sugeriu o desvio padrão, menor no grupo dos ingressantes do que no dos concluintes.
As notas máximas obtidas por concluintes ( 92,8 ) e ingressantes ( 92,3 ) foram diferentes.
As distribuições de notas para cada grupo se apresentaram com uma inclinação para notas menores, como evidenciados pelos coeficientes de assimetria positivos.
A situação de assimetria é confirmada pelo Gráfico 3.5 que proporciona uma comparação do desempenho entre concluintes e ingressantes em relação ao Componente Específico da prova.
Claramente a distribuição das notas médias dos ingressantes, à esquerda daquela dos concluintes, evidenciou o percentual mais elevado de notas médias nas faixas mais baixas para este grupo.
As medianas das distribuições das notas dos ingressantes e dos concluintes foram 35,6 e 38,1, respectivamente, enfatizando o melhor dos concluintes nesta parte da prova.
Assim como os Gráficos 3.2 e 3.4, o Gráfico 3.6 apresenta uma comparação dos resultados em relação às Regiões, Categoria Administrativa e Organização Acadêmica, agora levando em conta o desempenho dos ingressantes e dos concluintes no Componente Específico da prova.
Tendo como foco as notas médias dos estudantes, em determinada Região, observa -se que, entre os concluintes, as mais elevadas foram encontradas nas Regiões Norte ( 38,74 ), Sudeste ( 40,70 ) e Centro-Oeste ( 40,07 ), valores acima da média nacional para esta parte ( 38,67 ).
Entre os ingressantes, as maiores notas ocorreram nas Regiões Norte ( 37,03 ) e Sudeste ( 37,27 ), todas acima da média nacional ( 36,50 ).
Novamente, para a análise segundo a Categoria Administrativa, todas as entidades foram representadas.
Observa -se comportamento diverso àquele da Formação Geral, as médias mais elevadas entre os concluintes estão nas instituições federais ( 45,85 ), enquanto que na parte relativa à Formação Geral estavam nas estaduais.
As médias mais elevadas entre os ingressantes ( 42,79 ) ficam nas instituições estaduais, sendo que o mesmo ocorreu para a parte relativa à Formação Geral e ao Componente Específico.
Quanto à Organização Acadêmica, o comentário é o mesmo já feito para a nota bruta e referente à Formação Geral : somente Universidades, Centros Universitários, Faculdades e Institutos Federais de Educação, Ciência e Tecnologia são representadas neste curso.
Considerando -se o tipo de Organização Acadêmica, entre os concluintes, Instituto Federal de Educação, Ciência e Tecnologia apresentaram os maiores valores, seguidas dos Centros Universitários e das Faculdades.
E entre os ingressantes, Centros Universitários apresentaram os maiores valores, seguidos das Faculdades e das Universidades. 3.2 Análise das Questões Objetivas 3.2.1 Formação Geral A Tabela 3.4 apresenta as estatísticas básicas relativas às oito questões objetivas da Formação Geral dos estudantes.
Adiantamos que a primeira questão foi desconsiderada pelo critério da correlação ponto-bisserial.
Como pode ser observado na Tabela 3.4, os estudantes concluintes obtiveram desempenho superior, com nota média de 50,27, em relação aos estudantes ingressantes cuja nota média foi 48,85.
A variabilidade entre os concluintes foi um pouco menor do que entre os ingressantes, como indicaram os desvios-padrão de 21,8, para o grupo de ingressantes e 21,5 para o grupo dos concluintes.
A amplitude total das notas foi a mesma para ambos os grupos : nota mínima de 0,00 e máxima de 100,0.
A mediana foi superior para os concluintes : 57,1.
A assimetria, medida pelo coeficiente de assimetria ( skewness ) foi nula indicando simetria para os grupos.
A Tabela 3.5 apresenta a classificação das questões objetivas da Formação Geral, segundo o índice de facilidade ( questões em vermelho foram desconsideradas pelo critério ponto-bisserial ).
Como as questões foram comuns a todas as áreas, a tabela apresenta a classificação considerando -se os concluintes de todas as áreas e aqueles do curso de TECNOLOGIA EM PROCESSOS GERENCIAIS.
De acordo com os índices obtidos, as questões objetivas da prova foram assim avaliadas para a área de TECNOLOGIA EM PROCESSOS GERENCIAIS : das oito questões, duas tiveram entre 60 \%, inclusive, e 85 \% de acertos e foram classificadas como fáceis ( somente uma para o total de concluintes ) ; outras duas foram consideradas medianas, alcançando entre 40 \%, inclusive, e 60 \% de acertos ( três para o total de concluintes ) ; três obtiveram entre 15 \%, inclusive, e 40 \% de acertos, sendo consideradas difíceis ( três para o total de concluintes ).
Por fim, uma foi classificada como muito difícil tento obtido percentual menor do que 15 \% de acertos ( uma para o total de concluintes ), coincidentemente a questão desconsiderada pelo critério ponto-bisserial.
Nenhuma questão foi classificada na categoria muito fácil.
A Tabela 3.6 descreve os resultados das análises das questões objetivas relativas à Formação Geral, segundo o poder de discriminação, utilizando -se para tal o índice de discriminação ponto-bisserial.
Como as questões foram comuns a todas as áreas, a tabela apresenta a classificação considerando -se os concluintes de todas as áreas e aqueles do curso de TECNOLOGIA EM PROCESSOS GERENCIAIS.
A desconsideração da questão segundo este critério considerará o universo de todos os concluintes.
Nenhuma das oito questões apresentou índices iguais ou acima de 0,40 ( nenhuma para a população de concluintes como um todo ), sendo, portanto, classificadas como muito boas nesta característica, para este grupo de estudantes, uma foi classificada como boa com índice entre 0,30, inclusive, e 0,40 ( três para os concluintes como um todo ) e cinco foram classificadas como médias, com índice entre 0,20, inclusive, e 0,30 ( quatro para os concluintes como um todo ).
Duas questões foram classificadas como fraca tanto para a população de concluintes como um todo para a de concluintes de TECNOLOGIA EM PROCESSOS GERENCIAIS e foi desconsiderada pelo critério.
Na sequência, no Quadro 3.1, estão especificadas, por questão, as habilidades aferidas pelas questões de Formação Geral da prova da TECNOLOGIA EM PROCESSOS GERENCIAIS, bem como os resultados da análise relativa aos índices de facilidade e de discriminação de cada questão.
O índice de facilidade teve uma variação considerável, oscilando entre 0,11 até 0,80 ( na população de concluintes o intervalo foi de 0,08 a 0,75 ) e o de discriminação ficou entre 0,00 até 0,36 ( na população de concluintes o intervalo foi de 0,02 a 0,39 ).
A questão 7 aferiu as habilidades de “ Analisar e criticar informações ”, “ Extrair conclusões por indução e/ou dedução ”, “ Estabelecer relações, comparações e contrastes em diferentes situações ”, “ Detectar contradições ”, “ Fazer escolhas valorativas avaliando consequências ”, “ Questionar a realidade ”, e “ Argumentar coerentemente ” foi a que apresentou, para a população de concluintes de TECNOLOGIA EM PROCESSOS GERENCIAIS, maior poder discriminatório com índice de 0,36.
A questão mais fácil, com 80 \% de acertos, foi a de número 8, que aferiu também, além das já listadas, a habilidade de “ Ler e interpretar textos ”.
De outra parte, a questão com menor índice de facilidade e de discriminação deste em Formação Geral foi a questão 1, que foi desconsiderada, e que deveria aferir as habilidades de “ Analisar informações ” ; “ Extrair conclusões ” e “ Propor soluções para situações-problema ”.
O Gráfico 3.7 mostra, a título de ilustração, o comportamento das respostas à questão 7 de Formação Geral para os estudantes concluintes da área de TECNOLOGIA EM PROCESSOS GERENCIAIS como função da nota geral na prova4.
Trata -se de uma questão considerada de dificuldade fácil para os concluintes de TECNOLOGIA EM PROCESSOS GERENCIAIS e de dificuldade média para o conjunto de concluintes de todas as áreas.
O poder de discriminação foi considerado bom tanto para os estudantes da área, como também, para o conjunto de todos os concluintes.
Neste Gráfico, cada uma das cinco curvas representou o percentual de respostas em determinada alternativa da questão, em função do escore total dos estudantes concluintes da área de TECNOLOGIA EM PROCESSOS GERENCIAIS.
A curva azul marinho corresponde à alternativa A, a correta para este esta questão.
Assim, observa -se que entre os estudantes com notas mais baixas – desempenho geral mais fraco na prova como um todo –, além da alternativa correta, houve a escolha de alternativas incorretas, especialmente a alternativa B – curva verde claro seguida da alternativa E – curva em roxo.
A proporção de estudantes que selecionaram a resposta correta A, aumentou com o desempenho deste grupo de estudantes chegando ao acerto de 100 \% dos estudantes com notas pelo menos 87,5.
Esta análise permitiu verificar como a questão discriminou entre os grupos de desempenho.
Os Gráficos relativos às demais questões de Formação Geral constam do Anexo I. 3.2.2 Componente Específico A Tabela 3.7 apresenta as estatísticas básicas em relação às questões objetivas do Componente Específico.
Como pode ser observado, os estudantes concluintes obtiveram desempenho superior, mas por uma pequena margem com uma média de 41,04 pontos, em relação aos estudantes ingressantes com uma média de 38,57 pontos.
Os desvios-padrão indicaram que a variabilidade entre os concluintes foi maior ( 17,0 ) do que entre os ingressantes ( 16,3 ).
Os coeficientes de assimetria para os dois grupos foram positivos indicando uma inclinação para a esquerda.
A assimetria foi igual para ambos os grupos.
A análise das questões objetivas do Componente Específico em TECNOLOGIA EM PROCESSOS GERENCIAIS, em relação ao grau de facilidade da prova ( ver Tabela 3.8 ), mostrou que, das 27 questões, quatro encontraram -se na categoria ' muito difícil ' ( abaixo de 15 \% de acertos ), doze são de nível ' difícil ' ( entre 15 \% e 40 \% de acertos ), oito estão na categoria ' médio ' ( entre 40 \% e 60 \% de acertos ) e duas teve a classificação ' fácil ' ( entre 60 \% e 85 \% de acertos ).
Nenhuma questão foi classificada como ' muito fácil ' Uma questão foi desconsiderada pela comissão.
Das 27 questões objetivas relativas ao Componente Específico ( ver Tabela 3.9 ), três obtiveram o índice de discriminação muito bom, quatro obtiveram o índice de discriminação bom e cinco obtiveram nível de discriminação médio.
O índice fraco que eliminaria a questão foi alcançado por quatorze questões.
Uma questão foi cancelada pela Comissão Na sequência, no Quadro 3.2 estão especificadas, por questão, as habilidades aferidas pelas questões do Componente Específico da prova, bem como os resultados da análise relativa aos índices de facilidade e de discriminação de cada questão.
O Gráfico 3.8 mostra, a título de ilustração, o comportamento da questão 11 do Componente Específico para os estudantes concluintes da área de TECNOLOGIA EM PROCESSOS GERENCIAIS como função da nota geral na prova5.
Trata -se de uma questão considerada de dificuldade média para os concluintes de TECNOLOGIA EM PROCESSOS GERENCIAIS.
O poder de discriminação foi também médio para os estudantes da área.
Neste Gráfico, cada uma das cinco curvas, representa o percentual de respostas em determinada alternativa da questão, em função da nota total dos estudantes concluintes de TECNOLOGIA EM PROCESSOS GERENCIAIS.
A curva em roxo corresponde à alternativa E, a correta para esta questão.
Assim, observou -se que entre os estudantes com escores mais baixos – desempenho geral mais fraco na prova como um todo –, houve a escolha de alternativas incorretas como a alternativa B – curva verde claro ou a alternativa C – curva em vermelho.
A proporção de estudantes que selecionaram a resposta correta E aumentou com o desempenho deste grupo de estudantes chegando ao acerto de 100 \% dos estudantes com notas pelo menos 87,5.
Esta análise permitiu verificar como a questão discriminou entre os grupos de desempenho.
Os Gráficos relativos às demais questões do Componente Específico constam do Anexo I. 3.3 Análise das Questões Discursivas 3.3.1 Formação Geral A Tabela 3.10 apresenta as estatísticas básicas da prova por grupo de estudantes.
A análise dos resultados de desempenho dos estudantes nas questões discursivas que são apresentados na tabela 3.10 e no gráfico 3.9 mostrou que as notas foram mais baixas no conjunto dessas questões que no das questões objetivas.
O total de estudantes inscritos para a prova foi de 29.633.
Destes, 64,0 \% compareceram, sendo que a presença foi menor entre concluintes ( 62,6 \% ) do que entre ingressantes ( 65,0 \% ).
A média geral da prova foi de 28,54, sendo que os ingressantes obtiveram média mais baixa, de 27,71, que os concluintes, cuja média foi 29,83 O desvio-padrão geral foi de 33,6, sendo o do grupo de ingressantes menor ( 33,5 ), indicando que, em contraposição ao grupo dos concluintes ( 33,7 ), aqueles possuíram uma distribuição mais homogênea nas notas.
A nota máxima foi 100,0, atingida por concluintes e ingressantes.
Os coeficientes de assimetria para os dois grupos considerados, o de concluintes e o de ingressantes, bem como para a população como um todo são todos positivos, indicando uma maior concentração à esquerda da distribuição.
Para todos os grupos um pouco mais de metade dos estudantes presentes obtiveram nota zero nas questões discursivas da Formação Geral.
Das 9.859 provas com nota zero, 446 ( 4,5 \% ) continham um protesto, 2.223 ( 22,5 \% ) estavam em branco e 7.190 ( 72,9 \% ) foram respondidas, porém de forma inadequada.
Entre os concluintes, a maioria ( 72,7 \% ) dentre aqueles com nota zero estavam com as provas respondidas de forma inadequada.
Entre os ingressantes 73,0 \% das notas nulas decorreram de correção.
Uma comparação entre o desempenho de ingressantes e concluintes pode ser realizada a partir dos dados apresentados no Gráfico 3.9.
As notas dos concluintes, de um modo geral, foram ligeiramente mais elevadas do que as dos ingressantes, cuja distribuição de notas está à esquerda dos primeiros.
A ideia de melhor desempenho foi reforçada pela proporção de notas nulas : 50,4 \% entre os concluintes e 53,7 \% entre os ingressantes.
Na sequência, discorrer -se-á, especificamente, sobre uma das questões discursivas de Formação Geral, a número 10, já que a número 9 foi desconsiderada.
Em especial, serão analisados os conteúdos englobados na questão e os resultados obtidos pelos estudantes.
Ao final, são apresentados os comentários da banca corretora a respeito do observado nas respostas dos estudantes, suas impressões e conclusões.
Cumpre esclarecer que, tendo em vista que as questões discursivas de Formação Geral foram padronizadas, ou seja, constam de todas as provas, os comentários da banca foram os mesmos para todas as áreas acadêmicas, eis que direcionados a todos os estudantes que participaram do ENADE/2009. 3.3.1.1 Análise da Questão Discursiva 9 de Formação Geral Desconsiderada 3.3.1.2 Análise da Questão Discursiva 10 de Formação Geral A Questão Discursiva nº. 10 buscou avaliar as habilidades : • Ler e interpretar textos ; • Analisar e criticar informações ; • Extrair conclusões por indução e/ou dedução ; • Estabelecer relações, comparações e contrastes em diferentes situações ; • Detectar contradições ; • Fazer escolhas valorativas avaliando consequências ; • Questionar a realidade ; • Argumentar coerentemente.
O Quadro 3.3 apresenta as habilidades/competências, temas predominantes e estatísticas básicas da questão 10 de Formação Geral por grupo de estudantes.
Cumpre notar que como a questão 9 foi desconsiderada, a nota referente às questões discursivas da Formação Geral é idêntica à nota da questão 10.
A média geral desta questão foi de 28,54, sendo que os ingressantes obtiveram média mais baixa, de 27,71, que os concluintes, cuja média foi 29,83.
O desvio-padrão geral foi de 33,6, sendo que o do grupo de ingressantes menor ( 33,5 ), em relação ao grupo dos concluintes ( 33,7 ), aqueles possuíram uma distribuição mais homogênea nas notas.
A nota máxima foi 100,0, alcançada por concluintes e ingressantes.
Os coeficientes de assimetria para os dois grupos considerados, o de concluintes e o de ingressantes, bem como para a população como um todo foram todos positivos, indicando uma maior concentração à esquerda da distribuição.
Para todos os grupos, um pouco mais de metade dos estudantes presentes tiveram nota zero nesta questão discursiva de Formação Geral.
Das 9.859 provas com nota zero, 446 ( 4,5 \% ) continham um protesto, 2.223 ( 22,5 \% ) estavam em branco e 7.190 ( 72,9 \% ) foram respondidas, porém de forma inadequada.
Entre os concluintes a maioria ( 72,7 \% ) dentre aqueles com nota zero estavam com as provas respondidas de forma inadequada.
Entre os ingressantes 73,0 \% das notas nulas ( zero ) decorreram do desempenho do próprio estudante verificadas na correção.
Uma comparação entre o desempenho de ingressantes e concluintes para esta questão foi exatamente igual à já realizada para a nota discursiva de Formação Geral a partir dos dados apresentados no Gráfico 3.9. 3.3.1.3 Comentários sobre as respostas à Questão Discursiva 10 Apresentação A questão 10 de Formação Geral apresentou um questionamento capaz de analisar a habilidade do estudante em reconhecer o que é um direito social, optar por um direito social vigente que possa diminuir as desigualdades de renda familiar no Brasil e discutir os argumentos em favor do direito social escolhido como suporte à resposta.
Para análise das respostas apresentadas pelos estudantes, houve uma leitura inicial bastante criteriosa.
Foram considerados para efeito de pontuação, os seguintes argumentos : • Acesso à educação pública, gratuita e de qualidade, garantindo ao cidadão melhores condições de vida ; • Assistência à saúde, em seu contexto mais amplo, garantindo os cuidados essenciais de saúde, baseados em técnicas apropriadas ao alcance de todas as pessoas, contando com a participação da população ; • Permanência do estudante na escola, em todos os níveis escolares ( da educação infantil a educação superior, possibilitando a qualificação profissional ), o que garante melhores condições de trabalho e remuneração.
Considerou -se a qualificação e o desempenho docente, desenvolvendo nos estudantes melhores expectativas de aprendizagem e interesse pelo conteúdo da escola.
Os professores desenvolvendo metodologias adequadas e conhecimento pleno quanto a questões de aprendizagem ; • Condições dignas de trabalho, com remuneração que garanta qualidade de vida.
Foram consideradas as condições de trabalho, a preparação profissional para o trabalho e os cursos de aperfeiçoamento, oportunizando momentos de aprendizagem e desenvolvimento, a partir de situações desencadeadoras ; • Ser proprietário do imóvel que reside, reduzindo os gastos com aluguel, garantindo melhores condições de vida, equilíbrio financeiro e construção do bem-estar familiar.
Desempenho dos estudantes na questão É importante salientar que havia respostas desenvolveram adequadamente o questionamento apresentado, com fundamentação baseada em argumentos coerentes.
Contudo, em um número considerável de respostas, havia apenas um argumento e não dois, conforme fora solicitado.
Muitos estudantes informaram sobre problemas da educação, de um modo geral.
Algumas observações mencionadas pelos corretores durante o processo de correção foram : • há preocupação com uma educação de qualidade, com sugestões de cursos profissionalizantes e métodos preparatórios de aulas. • a educação é revelada como de baixa qualidade, exaltando -se treinamentos e cursos de atualização para a classe docente. • a saúde como essência para a vida.
O SUS é colocado como parte integrante de direito social, garantindo às pessoas fácil acesso. • oportunidade de emprego como forma de ascensão social, direito, cidadania. • os políticos brasileiros como causadores do despreparo educacional.
As respostas analisadas apresentaram norma culta, coesão e coerência razoáveis, que garantiram a compreensão da mensagem desenvolvida pelos estudantes.
Os erros ortográficos foram poucos e não comprometeram a discussão sobre o assunto e sobre os pontos de vista abordados.
Quanto à coerência, houve um desvio quanto ao objetivo da questão, já que alguns estudantes, através dos seus argumentos retrataram reivindicações, reclamações e situações de problemas enfrentados pelos brasileiros, citando golpes políticos.
Porém, mesmo quando se tratava desta avaliação, os estudantes foram autênticos através das explicações.
Houve, também, erros quanto à concordância verbal e nominal.
Houve também, dificuldade em utilizar os conectores corretamente.
Alguns estudantes, por falta de ideias, apresentaram em suas respostas redundância e prolixidade.
O desempenho foi considerado satisfatório, pois a maioria dos estudantes soube responder ao questionamento.
Houve domínio do conteúdo, pois as ideias, os argumentos e a originalidade corresponderam ao esperado. 3.3.2 Componente Específico Nas questões discursivas do Componente Específico – questões de números 38, 39 e 40 –, a diferença de desempenho entre os grupos não foi muito acentuada, inclusive menor do que para as questões discursivas de Formação Geral.
A Tabela 3.11 contém as estatísticas da média destas questões.
Observa -se que no grupo dos concluintes, a nota média alcançada foi 25,64, enquanto que no grupo dos ingressantes foi de 25,11, ambas menores do que as obtidas nas questões discursivas de Formação Geral, que foram de 27,71 para os ingressantes e 29,83 para os concluintes ( Tabela 3.10 ).
As médias foram também menores dos que as obtidas para as questões objetivas do Componente Específico ( Tabela 3.7 ) : respectivamente 38,57 e 41,04 para ingressantes e concluintes.
A análise de cada uma destas questões será feita nas subseções seguintes.
Os valores positivos para a medida de assimetria indicaram uma distribuição com cauda para a esquerda.
Dos 29.633 presentes, 24,5 \% obtiveram nota zero.
A tabela 3.12 apresenta um detalhamento da situação destes estudantes.
A proporção dos que não fizeram nenhuma das 3 questões foi de 37,1 \%.
Esta proporção foi maior entre os ingressantes ( 40,0 \% ) do que entre concluintes ( 32,6 \% ).
A proporção de provas com todas as questões respondidas correspondeu a 53,0 \% das provas.
Esta proporção foi maior entre os concluintes ( 58,5 \% ) do que entre os ingressantes ( 49,4 \% ).
Houve apenas 1 prova com as 3 questões com protesto.
O Gráfico 3.10 apresenta a distribuição das notas nas questões discursivas do Componente Específico por grupo de estudantes.
A moda se situou no primeiro intervalo, aquele com valores entre 0 e 10 incluindo os extremos do intervalo.
A impressão de assimetria informada pelos coeficientes de assimetria positivos ( Tabela 3.11 ) foi reforçada no gráfico pela concentração no primeiro intervalo. 3.3.2.1 Análise da Questão Discursiva 38 do Componente Específico A questão 38, cujas habilidades aferidas a serem aferidas encontram -se descritas no Quadro 3.4, que ainda contém outras informações, seguido da Tabela 3.13 e do Gráfico 3.11, foi aquela em que os grupos de estudantes observados obtiveram o pior desempenho entre as questões discursivas do Componente Específico.
A nota média do grupo dos ingressantes e a dos concluintes foi igual a 22,93.
Nesta questão ocorre também o pior percentual de notas nulas : nos dois grupos foi de 43,3 \%, quase igual para os ingressantes ( 43,2 \% ) e concluintes ( 43,4 \% ).
Quase 40 \% dos concluintes e dos ingressantes que tiraram zero deixaram a questão em branco ( ver Tabela 3.13 ).
Um grupo pequeno ( 0,9 \% ) utilizou a questão para protestar.
Apesar do baixo desempenho, em ambos os grupos a nota máxima de 100,0 pontos foi alcançada.
Os coeficientes de assimetria positivos indicam uma inclinação para à esquerda e cauda para a direita.
O Gráfico 3.11 apresenta a distribuição das notas na questão 38 discursiva do Componente Específico por grupo de estudantes.
A moda se situou no primeiro intervalo, aquele com valores entre 0 e 10 incluindo os extremos do intervalo.
A impressão de assimetria informada pelos coeficientes de assimetria positivos ( Quadro 3.4 ) foi reforçada no gráfico pela concentração no primeiro intervalo 3.3.2.2 Comentários sobre as respostas à Questão Discursiva 38 Essa questão teve como objetivo identificar e explicar corretamente as ações relacionadas as quatro funções administrativas, verificar o entendimento dos estudantes em relação às atividades das funções gerenciais de planejamento, organização, direção e controle de um estudo de reestruturação organizacional.
Possui nível de dificuldade mediano.
Quanto ao desempenho dos estudantes, pode -se afirmar que foi razoável no que se refere à apresentação da resposta em relação ao domínio do conteúdo.
A falta de compreensão do foco da questão levou um grande número de estudantes apenas a repetir o enunciado, o que demonstrou também, a falta de clareza e a concepção de que o ato de copiar não induziria o corretor.
Neste caso, entendemos que ocorreu um fraco desempenho.
No que se refere ao texto escrito e linguagem empregada, entendemos que a grande maioria dos estudantes necessitou de reforço para interpretar e redigir, daí a falta de entendimento ao enunciado da questão.
Assim, evidencia -se que alguns estudantes confundiram organização e planejamento e não classificaram as medidas, demonstrando desconhecimento. 3.3.2.3 Análise da Questão Discursiva 39 do Componente Específico A questão 39, cujas habilidades aferidas encontram -se descritas no Quadro 3.5, que ainda contém outras informações, seguido da Tabela 3.14 e do Gráfico 3.12, foi aquela em que os grupos de estudantes observados obtiveram um desempenho intermediário entre as questões discursivas do Componente Específico.
A nota média do grupo dos ingressantes foi de 25,54, e a dos concluintes 26,33.
Nesta questão ocorre também o menor percentual de notas nulas : nos dois grupos foi de 36,1 \% entre os ingressantes e 35,1 \% entre os concluintes.
Das provas com notas nulas nesta questão, quase 35 \% foi deixada em branco : 36,4 \% entre os ingressantes e 28,9 \% entre os concluintes ( ver Tabela 3.14 ).
Entre os com nota nula, uma pequena minoria ( 0,7 \% ) utilizou a questão para protestar.
Apesar do baixo desempenho, em ambos os grupos a nota máxima de 100,0 pontos foi alcançada.
Os coeficientes de assimetria positivos indicaram uma inclinação para à esquerda e cauda para a direita.
O Gráfico 3.12 apresenta a distribuição das notas na questão 39 discursiva do Componente Específico por grupo de estudantes.
A moda se situou no primeiro intervalo, aquele com valores entre 0 e 10 incluindo os extremos do intervalo.
A impressão de assimetria informada pelos coeficientes de assimetria positivos ( Quadro 3.5 ) foi reforçada no gráfico pela concentração no primeiro intervalo. 3.3.2.4 Comentários sobre as respostas à Questão Discursiva 39 A questão em pauta teve como objetivo explicitar todas as características de um empreendedor ( oportunidade, ideia, implementação e resultado ) e explicar corretamente o conceito de mercado para produtos inovadores.
Portanto, foi uma questão com exigência de duas respostas A e B que se completavam, facilitando o entendimento.
O grau de dificuldade da questão foi mediano e os estudantes tiveram um desempenho fraco, pois não dominaram o conteúdo.
O que mais surpreendeu é que houve muita dificuldade na articulação de conceitos, mesmo se tratando de um estudo de caso.
A grande maioria dos estudantes apresentou apenas uma resposta, o que demonstrou falta de compreensão do enunciado.
Usaram poucas palavras e muitas delas, escritas incorretamente.
Concluindo, pode -se afirmar que alguns estudantes pareceram ter dificuldades na resposta, pela falta de leitura do questionamento.
Aparentemente, os estudantes leram o estudo de caso e passaram a emitir suas opiniões, sem dedicar atenção à pergunta da questão.
Muitos estudantes acertaram o item A da questão, mas não responderam o item B da mesma.
Alguns estudantes, tiveram dificuldade de articular os conceitos com a situação prática do estudo de caso, tratando a empresa inovadora do Sr.
Peter, como sendo uma empresa produtora de commodities no mercado.
Outro ponto interessante foi a economia de palavras para responder a questão.
Os estudantes pareciam estar desmotivados em demonstrar seus conhecimentos, respondendo de maneira sucinta conceitos complexos, como por exemplo, as características de um empreendedor. 3.3.2.5 Análise da Questão Discursiva 40 do Componente Específico A questão 40, cujas habilidades aferidas encontram -se descritas no Quadro 3.6, que ainda contém outras informações, seguido da Tabela 3.15 e do Gráfico 3.13.
A questão foi aquela em que os grupos de estudantes observados obtiveram o melhor desempenho entre as questões discursivas do Componente Específico.
A nota média do grupo dos ingressantes foi de 26,93, e a dos concluintes 27,75.
O percentual de notas nulas foi de 39,7 \% entre os ingressantes e 38,7 \% entre os concluintes.
Aqui a maioria das provas com nota zero decorreu de provas mal respondidas : 61,8 \% entre os ingressantes e 68,3 \% entre os concluintes ( ver Tabela 3.15 ).
Uma pequena minoria ( 0,6 \% ), dentre aqueles com nota nula, utilizou a questão para protestar.
Em ambos os grupos a nota máxima de 100,0 pontos foi alcançada.
Os coeficientes de assimetria positivos indicaram uma inclinação para a esquerda.
O Gráfico 3.13 apresenta a distribuição das notas na questão 40 discursiva do Componente Específico por grupo de estudantes.
A moda se situou no primeiro intervalo, aquele com valores entre 0 e 10 incluindo os extremos do intervalo.
A impressão de assimetria informada pelos coeficientes de assimetria positivos ( Quadro 3.6 ) foi reforçada no gráfico pela concentração no primeiro intervalo. 3.3.2.6 Comentários sobre as respostas à Questão Discursiva 40 Esta é uma questão de nível de dificuldade fácil que tem como objetivo identificar as vantagens e as desvantagens de forma correta e completa para os dois planos de ação.
Sendo uma questão direta, em função de uma situação problema muito próxima dos estudantes desta área, o desempenho foi razoável.
O critério adotado para correção foi claro e todos os corretores dimensionaram a questão de forma a aproveitar satisfatoriamente os conceitos emitidos pelos estudantes.
Assim, sendo, muitos estudantes responderam somente vantagens e desvantagens do primeiro plano, omitindo as respostas para o segundo plano.
Isso demonstra falta de compreensão do enunciado e articulação pobre nas respostas.
Finalizado, pode -se afirmar que em se tratando de questão que abordava o assunto mudança organizacional, solicitando diretamente as vantagens e desvantagens de se contratar um novo grupo de gerentes e também listar as vantagens e as desvantagens de conservar o time de gerentes atuais e treiná -los de forma intensiva para realizar a mudança de cultura, não justificou o desempenho fraco dos estudantes.

% matched lemmas: 
\end{textsample}
